\chapter{Conclusion}
\label{ch:conclusion}

\section{Summary}

This study asked whether a deep fall detector trained on healthy volunteers can recognise freezing of
gait in people with Parkinson's disease, and why it cannot. With the confounders removed by
construction, the answers to the four research questions are as follows.

\textbf{RQ1.} On unseen SisFall subjects, CNN-1D, LSTM, GRU and patch Transformer are all excellent
(AUROC 0.988--0.999, 0.7--1.5 false alarms per hour). The recurrent models are marginally ahead,
but no ranking is supported by the subject-level uncertainty. With impact-centred labels, simulated fall
detection is essentially a solved problem, so it is a poor benchmark for choosing an architecture.

\textbf{RQ2.} Zero-shot, the detectors rank FoG only slightly above walking (AUROC 0.49--0.63; 0.65 for the
Transformer seed ensemble, the only interval that excludes chance), and at their validated operating
point they detect essentially no freezing episode while raising 12--23 false alarms per hour of
Parkinsonian walking. A single-line freeze index outperforms all of them.

\textbf{RQ3.} The gap is mostly concept shift. Label-free normalisation reduces the distance between
domains without making them indistinguishable and without improving transfer, and the direction that
the networks use to flag falls is nearly orthogonal to the one that separates FoG from walking.

\textbf{RQ4.} When patients are labelled, fall pre-training is a worse initialisation than random
weights. Trained from scratch on six patients with a trunk sensor, the models reach an AUROC of about
0.78, with the CNN-1D and the Transformer the most data-efficient.

The practical lesson is simple: a monitoring system for Parkinson's disease cannot inherit its notion
of ``abnormal'' from simulated falls. A model that is near-perfect on the source domain can be blind
to the target event, and nothing measured on the source domain warns about it.

\section{Perspectives}

\paragraph{Event-agnostic pre-training.} Fall supervision teaches impact detectors. Self-supervised
objectives on large unlabelled wearable corpora (contrastive or masked reconstruction,
\cite{eldele2021tstcc}) learn representations of movement itself, including the rhythmic structure that
FoG disrupts, and are a better candidate for transfer.

\paragraph{Physics as an inductive bias.} The freeze index beat every network without training. Feeding
spectral features in the locomotion and freeze bands to the models, or constraining the first layers to
behave like band-pass filters, could combine the robustness of the physical prior with the capacity of
learned models.

\paragraph{Personalisation.} Per-patient AUROC varied from 0.26 to 0.85, and zero-shot rankings even
reversed between patients. Few-shot adaptation from a handful of annotated episodes per patient, or
test-time adaptation of normalisation layers \parencite{li2017adabn}, is a natural next step.

\paragraph{Better targets.} Daphnet is small and was recorded in a laboratory protocol designed to
provoke freezing. Larger home-recorded cohorts with both FoG and real falls, and sensors at the shank as
well as the trunk, would make it possible to test whether the conclusions hold outside this setting.
